% Part: incompleteness
% Chapter: theories-computability
% Section: inseparability

\documentclass[../../../include/open-logic-section]{subfiles}

\begin{document}

\olfileid{inc}{tcp}{ins}

\olsection{$\Th{Q}$ मध्ये सिद्ध होणारी आणि ज्यांचे निषेध सिद्ध होतात
  अशा वाक्यांचे संच संगणनक्षमपणे अविलगनीय आहेत}


$\Th{\bar Q}$ हा ज्या वाक्यांचे {\em निषेध} $\Th{Q}$ मध्ये
सिद्ध होतात त्यांचा संच असू द्या, म्हणजे
$\Th{\bar Q} = \Setabs{!A}{\Th{Q} \Proves
  \lnot !A}$. विभक्त संच $A$ आणि $B$ यांना
संगणनक्षमपणे अविलगनीय म्हणतात, जर असा कोणताही
संगणनक्षम संच $C$ नसेल की $A \subseteq C$ आणि
$B \subseteq \Complement{C}$.

\begin{lem}
$\Th{Q}$ आणि $\Th{\bar Q}$ हे संगणनक्षमपणे अविलगनीय आहेत.
\end{lem}

\begin{proof}
$C$ हा असा संगणनक्षम संच आहे असे समजू की
$\Th{Q} \subseteq C$ आणि
$\Th{\bar Q} \subseteq \Complement{C}$.
$R(x,y)$ हा पुढील संबंध असू द्या:
\begin{quote}
$x$ हा $!D(u)$ या सूत्राचा संकेतांक आहे आणि
$!D(\num y)$ हे $C$ मध्ये आहे.
\end{quote}
सूत्र-संकेतांक ओळखणे, संख्यांकाचे प्रतिस्थापन करणे आणि
या संचातील सदस्यत्व तपासणे संगणनक्षम असल्यामुळे $R(x,y)$
संगणनक्षम आहे. आता तो सार्वत्रिक संगणनक्षम संबंध आहे असे दाखवू;
त्यातून व्याघात मिळेल.

$S(y)$ हा संगणनक्षम संबंध असून त्याचे $\Th{Q}$ मध्ये
$!D_S(u)$ या सूत्राने निरूपण होते असे समजू. मग
\begin{eqnarray*}
S(n) & \lif & \Th{Q} \Proves !D_S(\num n) \\
& \lif & !D_S(\num n) \in C
\end{eqnarray*}
आणि
\begin{eqnarray*}
\lnot S(n) & \lif & \Th{Q} \Proves \lnot !D_S(\num n) \\
& \lif & !D_S(\num n) \in \Th{\bar Q} \\
& \lif & !D_S(\num n) \not\in C
\end{eqnarray*}
असे मिळते. म्हणून $S(y)$ हे
$R(\#(!D_S(u)),y)$ शी सममूल्य आहे.
\end{proof}

\end{document}
